\chapter{Background and Context}


Chess remains to be the most popular board game in the world, with millions of players online on platforms such as \href{https://www.chess.com/}{\textcolor{blue}{chess.com}} and \href{https://lichess.org/}{\textcolor{blue}{lichess.org}}. The strategy and problem-solving skills required to solve certain positions creates a high skill ceiling, one which only few in the world have been able to reach. However, through recent technological advancements in processing power and memory storage, chess engines have been created that rival the ability of the best players, some even going beyond their skill level entirely. This is evident in the currently leading open-source chess engine \emph{Stockfish}, with a rating of approximately 3530 according to \href{https://www.chess.com/terms/stockfish-chess-engine}{\textcolor{blue}{chess.com}}, leagues beyond the highest rating achieved by a human ever \href{https://www.chess.com/article/view/best-chess-players-of-all-time#magnus}{\textcolor{blue}{(Magnus Carlsen at 2889)}}. \\
\\
Stockfish was able to achieve such a level through its recursive approach in "solving" chess, implementing the alpha-beta search algorithm at a depth higher than most other engines. The reason it was able to reach this depth was its pruning of branches that were less likely to produce good moves, which were assumed to be present later in the list of ordered moves. This approach was taken by many chess engines, but the aggressive \href{https://www.chessprogramming.org/Late_Move_Reductions}{\textcolor{blue}{LMR}} of Stockfish set it out amongst its competitors. \\
\\
This project aims to replicate the implementation of Stockfish by using the same alpha-beta search algorithm, but pair it with a classifier to prune branches deemed disadvantageous. This classifier will be built using the unsupervised machine learning algorithm, K-Means clustering.